% ============================================================
\chapter{Variantes de GNN (GAT, GIN, GraphSAGE)}
\label{ch:variantes}
% ============================================================

Le GCN de Kipf et Welling (2017) a ouvert la voie, mais ses limitations sont vite apparues~: agr\'egation isotrope, expressivit\'e plafonn\'ee au test de Weisfeiler-Leman d'ordre 1, incapacit\'e \`a traiter des graphes de grande taille de mani\`ere inductive. En r\'eponse, une vague d'architectures alternatives a d\'eferl\'e~: les \emph{Graph Attention Networks} (GAT) introduisent des m\'ecanismes d'attention pour pond\'erer les voisins~; les \emph{Graph Isomorphism Networks} (GIN) maximisent l'expressivit\'e th\'eorique~; \emph{GraphSAGE} permet l'apprentissage inductif par \'echantillonnage de voisinages. Ce chapitre explore ces variantes et les compromis qu'elles incarnent.

\section{Au-del\`a du GCN~: motivations}

Le GCN de Kipf et Welling pr\'esente plusieurs limitations~:
\begin{itemize}
    \item Agr\'egation \textbf{isotrope}~: tous les voisins sont trait\'es de la m\^eme
          mani\`ere (pond\'er\'es uniquement par le degr\'e).
    \item \textbf{Expressivit\'e limit\'ee}~: \'equivalence stricte au 1-WL.
    \item \textbf{Scalabilit\'e}~: n\'ecessite la matrice d'adjacence compl\`ete.
\end{itemize}

Ce chapitre pr\'esente trois variantes majeures qui r\'epondent \`a ces limitations.

\section{Graph Attention Networks (GAT)}

\subsection{Principe de l'attention sur graphes}

\begin{definition}[Couche GAT --- Veli\v{c}kovi\'c et al., 2018]
La couche GAT utilise un m\'ecanisme d'\textbf{attention} pour pond\'erer les
contributions des voisins~:
\[
    \mathbf{h}_v^{(\ell+1)} = \sigma\!\left(\sum_{u \in \mathcal{N}(v) \cup \{v\}}
    \alpha_{vu}^{(\ell)}\, \mathbf{W}^{(\ell)} \mathbf{h}_u^{(\ell)}\right)
\]
o\`u les coefficients d'attention $\alpha_{vu}$ sont calcul\'es par~:
\begin{align}
    e_{vu}^{(\ell)} &= \text{LeakyReLU}\!\left(\mathbf{a}^{\top}\!
    \left[\mathbf{W}\mathbf{h}_v^{(\ell)} \| \mathbf{W}\mathbf{h}_u^{(\ell)}\right]\right) \\
    \alpha_{vu}^{(\ell)} &= \frac{\exp(e_{vu}^{(\ell)})}{\sum_{w \in \mathcal{N}(v) \cup \{v\}}
    \exp(e_{vw}^{(\ell)})}
\end{align}
o\`u $\mathbf{a} \in \R^{2d'}$ est le vecteur d'attention appris et $\|$ d\'enote
la concat\'enation.
\end{definition}

\begin{remarque}[Attention multi-t\^ete]
Pour stabiliser l'apprentissage, GAT utilise $K$ t\^etes d'attention
ind\'ependantes~:
\[
    \mathbf{h}_v^{(\ell+1)} = \Big\|_{k=1}^{K} \sigma\!\left(\sum_{u \in \mathcal{N}(v)}
    \alpha_{vu}^{k} \mathbf{W}^k \mathbf{h}_u^{(\ell)}\right)
\]
La derni\`ere couche utilise typiquement la moyenne au lieu de la concat\'enation.
\end{remarque}

\begin{intuition}[title=\textbf{Pourquoi l'attention~?}]
L'attention permet au mod\`ele d'apprendre \textbf{quels voisins sont importants}
pour chaque n\oe ud, contrairement au GCN o\`u la pond\'eration est fix\'ee par
la structure du graphe. Cela est particuli\`erement utile quand~:
\begin{itemize}
    \item Les voisins ont des \textbf{pertinences variables}.
    \item Le graphe contient du \textbf{bruit} (ar\^etes non informatives).
    \item Les relations sont \textbf{asym\'etriques} (graphes dirig\'es).
\end{itemize}
\end{intuition}

\begin{codeblock}[title=\textbf{Impl\'ementation GAT avec PyTorch Geometric}]
\begin{minted}{python}
import torch
import torch.nn.functional as F
from torch_geometric.nn import GATConv
from torch_geometric.datasets import Planetoid

class GAT(torch.nn.Module):
    def __init__(self, in_channels, hidden_channels, out_channels,
                 heads=8, dropout=0.6):
        super().__init__()
        self.conv1 = GATConv(in_channels, hidden_channels,
                             heads=heads, dropout=dropout)
        self.conv2 = GATConv(hidden_channels * heads, out_channels,
                             heads=1, concat=False, dropout=dropout)
        self.dropout = dropout

    def forward(self, x, edge_index):
        x = F.dropout(x, p=self.dropout, training=self.training)
        x = F.elu(self.conv1(x, edge_index))
        x = F.dropout(x, p=self.dropout, training=self.training)
        x = self.conv2(x, edge_index)
        return x

dataset = Planetoid(root='/tmp/Cora', name='Cora')
data = dataset[0]
model = GAT(dataset.num_features, 8, dataset.num_classes, heads=8)
optimizer = torch.optim.Adam(model.parameters(), lr=0.005, weight_decay=5e-4)

for epoch in range(200):
    model.train()
    optimizer.zero_grad()
    out = model(data.x, data.edge_index)
    loss = F.cross_entropy(out[data.train_mask], data.y[data.train_mask])
    loss.backward()
    optimizer.step()

model.eval()
pred = model(data.x, data.edge_index).argmax(dim=1)
acc = (pred[data.test_mask] == data.y[data.test_mask]).float().mean()
print(f"GAT - Precision test: {acc:.4f}")
\end{minted}
\end{codeblock}

\subsection{GATv2~: attention dynamique}

\begin{definition}[GATv2 --- Brody et al., 2022]
GATv2 corrige une limitation de GAT o\`u l'attention est \textbf{statique}
(le classement des voisins ne d\'epend pas du n\oe ud requête). GATv2 utilise~:
\[
    e_{vu} = \mathbf{a}^\top \text{LeakyReLU}\!\left(\mathbf{W}\!\left[
    \mathbf{h}_v \| \mathbf{h}_u\right]\right)
\]
La non-lin\'earit\'e est appliqu\'ee \textbf{avant} le produit scalaire avec
$\mathbf{a}$, rendant l'attention v\'eritablement \textbf{dynamique}.
\end{definition}

\section{Graph Isomorphism Network (GIN)}

\subsection{Maximiser le pouvoir expressif}

\begin{theoreme}[GIN --- Xu et al., 2019]
Pour qu'un MPNN atteigne la puissance maximale du test 1-WL, les conditions
suivantes sont n\'ecessaires et suffisantes~:
\begin{enumerate}
    \item L'agr\'egation doit \^etre \textbf{injective} sur les multi-ensembles.
    \item La fonction de mise \`a jour doit \^etre \textbf{injective}.
\end{enumerate}
\end{theoreme}

\begin{definition}[Couche GIN]
La couche GIN r\'ealise~:
\[
    \mathbf{h}_v^{(\ell+1)} = \text{MLP}^{(\ell)}\!\left((1 + \epsilon^{(\ell)})
    \cdot \mathbf{h}_v^{(\ell)} + \sum_{u \in \mathcal{N}(v)}
    \mathbf{h}_u^{(\ell)}\right)
\]
o\`u $\epsilon$ est un param\`etre appris ou fix\'e.
\end{definition}

\begin{proposition}[Injectivit\'e de la somme]
Parmi les agr\'egations classiques (somme, moyenne, max), seule la \textbf{somme}
est injective sur les multi-ensembles~:
\begin{itemize}
    \item $\text{max}\{1, 1, 2\} = \text{max}\{1, 2\} = 2$ (perte d'information).
    \item $\text{mean}\{1, 1\} = \text{mean}\{1\} = 1$ (perte de cardinalit\'e).
    \item $\text{sum}\{1, 1, 2\} = 4 \neq \text{sum}\{1, 2\} = 3$ (injectif).
\end{itemize}
\end{proposition}

\begin{codeblock}[title=\textbf{Impl\'ementation GIN}]
\begin{minted}{python}
import torch
import torch.nn as nn
import torch.nn.functional as F
from torch_geometric.nn import GINConv, global_add_pool
from torch_geometric.datasets import TUDataset
from torch_geometric.loader import DataLoader

class GIN(nn.Module):
    def __init__(self, num_features, hidden_dim, num_classes, num_layers=5):
        super().__init__()
        self.convs = nn.ModuleList()
        self.bns = nn.ModuleList()
        for i in range(num_layers):
            in_dim = num_features if i == 0 else hidden_dim
            mlp = nn.Sequential(
                nn.Linear(in_dim, hidden_dim),
                nn.ReLU(),
                nn.Linear(hidden_dim, hidden_dim),
            )
            self.convs.append(GINConv(mlp, train_eps=True))
            self.bns.append(nn.BatchNorm1d(hidden_dim))
        self.classifier = nn.Linear(hidden_dim, num_classes)

    def forward(self, data):
        x, edge_index, batch = data.x, data.edge_index, data.batch
        for conv, bn in zip(self.convs, self.bns):
            x = F.relu(bn(conv(x, edge_index)))
        x = global_add_pool(x, batch)
        return self.classifier(x)

dataset = TUDataset(root='/tmp/MUTAG', name='MUTAG')
loader = DataLoader(dataset, batch_size=32, shuffle=True)
model = GIN(dataset.num_features, 64, dataset.num_classes)
print(f"Parametres: {sum(p.numel() for p in model.parameters()):,}")
\end{minted}
\end{codeblock}

\section{GraphSAGE~: apprentissage inductif}

\subsection{Motivation~: passage \`a l'\'echelle}

\begin{definition}[GraphSAGE --- Hamilton et al., 2017]
GraphSAGE (\textit{SAmple and aggreGatE}) utilise un \textbf{\'echantillonnage
de voisinage} pour permettre le passage \`a l'\'echelle~:
\[
    \mathbf{h}_v^{(\ell+1)} = \sigma\!\left(\mathbf{W}^{(\ell)} \cdot
    \text{CONCAT}\!\left(\mathbf{h}_v^{(\ell)},\;
    \text{AGG}(\{\mathbf{h}_u^{(\ell)} : u \in \mathcal{S}(v)\})\right)\right)
\]
o\`u $\mathcal{S}(v) \subseteq \mathcal{N}(v)$ est un sous-ensemble
\'echantillonn\'e de voisins (typiquement $\abs{\mathcal{S}(v)} \leq k$).
\end{definition}

\begin{proposition}[Avantages de GraphSAGE]
\begin{enumerate}
    \item \textbf{Inductif}~: peut g\'en\'eraliser \`a des n\oe uds non vus.
    \item \textbf{Scalable}~: complexit\'e contr\^ol\'ee par le nombre d'\'echantillons.
    \item \textbf{Mini-batch}~: compatible avec l'entra\^inement par mini-lots.
\end{enumerate}
\end{proposition}

\begin{algorithme}[title=\textbf{\'Echantillonnage de sous-graphes pour GraphSAGE}]
\begin{enumerate}
    \item Pour chaque n\oe ud cible $v$, \'echantillonner $k_1$ voisins directs.
    \item Pour chacun de ces voisins, \'echantillonner $k_2$ voisins (2-hop).
    \item Former le sous-graphe induit et propager les messages.
    \item La complexit\'e par batch est $\mathcal{O}(B \cdot k_1 \cdot k_2)$
          au lieu de $\mathcal{O}(n)$.
\end{enumerate}
\end{algorithme}

\begin{codeblock}[title=\textbf{GraphSAGE avec \'echantillonnage}]
\begin{minted}{python}
import torch
import torch.nn.functional as F
from torch_geometric.nn import SAGEConv
from torch_geometric.loader import NeighborLoader
from torch_geometric.datasets import Planetoid

dataset = Planetoid(root='/tmp/Cora', name='Cora')
data = dataset[0]

class GraphSAGE(torch.nn.Module):
    def __init__(self, in_channels, hidden_channels, out_channels):
        super().__init__()
        self.conv1 = SAGEConv(in_channels, hidden_channels)
        self.conv2 = SAGEConv(hidden_channels, out_channels)

    def forward(self, x, edge_index):
        x = F.relu(self.conv1(x, edge_index))
        x = F.dropout(x, p=0.5, training=self.training)
        x = self.conv2(x, edge_index)
        return x

# Echantillonnage de voisinage pour mini-batch
train_loader = NeighborLoader(
    data,
    num_neighbors=[25, 10],  # k1=25, k2=10
    batch_size=256,
    input_nodes=data.train_mask,
)

model = GraphSAGE(dataset.num_features, 64, dataset.num_classes)
optimizer = torch.optim.Adam(model.parameters(), lr=0.01)

for epoch in range(20):
    model.train()
    for batch in train_loader:
        optimizer.zero_grad()
        out = model(batch.x, batch.edge_index)[:batch.batch_size]
        loss = F.cross_entropy(out, batch.y[:batch.batch_size])
        loss.backward()
        optimizer.step()
\end{minted}
\end{codeblock}

\section{Comparaison des variantes}

\begin{center}
\begin{tabular}{lcccc}
\toprule
\textbf{Propri\'et\'e} & \textbf{GCN} & \textbf{GAT} & \textbf{GIN} & \textbf{SAGE} \\
\midrule
Attention & Non & Oui & Non & Non \\
Expressivit\'e 1-WL & $<$ & $<$ & $=$ & $<$ \\
Inductif & Non & Oui & Oui & Oui \\
Scalable (mini-batch) & Difficile & Mod\'er\'e & Mod\'er\'e & Oui \\
Attributs d'ar\^etes & Non & Oui (v2) & Non & Non \\
\bottomrule
\end{tabular}
\end{center}

\section{Architectures avanc\'ees}

\subsection{PNA --- Principal Neighbourhood Aggregation}

\begin{definition}[PNA --- Corso et al., 2020]
PNA combine \textbf{plusieurs agr\'egateurs} et \textbf{scalers}~:
\[
    \mathbf{h}_v^{(\ell+1)} = U\!\left(\mathbf{h}_v^{(\ell)},
    \underset{\substack{\oplus \in \{\text{sum, mean,}\\\text{max, std}\}}}
    {\Big\|}
    \underset{s \in \{1, \sqrt{d}, 1/\sqrt{d}\}}{\Big\|}
    s \cdot \bigoplus_{u \in \mathcal{N}(v)} M(\mathbf{h}_v^{(\ell)},
    \mathbf{h}_u^{(\ell)})\right)
\]
La concat\'enation de $4 \times 3 = 12$ canaux capture diff\'erents aspects
de la distribution du voisinage.
\end{definition}

\subsection{Connexions r\'esiduelles pour GNN profonds}

\begin{definition}[GCNII --- Chen et al., 2020]
GCNII ajoute une \textbf{connexion r\'esiduelle initiale} et un
\textbf{r\'etr\'ecissement d'identit\'e} pour \'eviter le sur-lissage~:
\[
    \mathbf{H}^{(\ell+1)} = \sigma\!\left(\left((1-\alpha_\ell)\hat{\mathbf{P}}
    \mathbf{H}^{(\ell)} + \alpha_\ell \mathbf{H}^{(0)}\right)
    \left((1-\beta_\ell)\mathbf{I} + \beta_\ell \mathbf{W}^{(\ell)}\right)\right)
\]
o\`u $\hat{\mathbf{P}} = \hat{\mathbf{D}}^{-1/2}\hat{\mathbf{A}}\hat{\mathbf{D}}^{-1/2}$
et $\alpha_\ell, \beta_\ell$ sont des hyperparam\`etres.
\end{definition}

\section*{Exercices}

\begin{exercice}[Comparaison exp\'erimentale]
Comparer GCN, GAT, GIN et GraphSAGE sur le dataset Cora~:
\begin{enumerate}
    \item Impl\'ementer les quatre mod\`eles avec le m\^eme nombre de param\`etres.
    \item Entra\^iner 200 \'epoques avec les m\^emes hyperparam\`etres.
    \item Reporter la pr\'ecision test et la courbe de perte.
    \item Analyser les diff\'erences de performance.
\end{enumerate}
\end{exercice}

\begin{exercice}[Visualisation de l'attention GAT]
\begin{enumerate}
    \item Entra\^iner un GAT sur Cora et extraire les coefficients d'attention $\alpha_{vu}$.
    \item Visualiser les poids d'attention sur un sous-graphe de 20 n\oe uds.
    \item Les coefficients d'attention sont-ils corr\'el\'es avec les labels des n\oe uds~?
\end{enumerate}
\end{exercice}

\begin{exercice}[Pouvoir expressif du GIN]
\begin{enumerate}
    \item Construire deux graphes que le GCN ne distingue pas mais que le GIN distingue.
    \item V\'erifier exp\'erimentalement.
    \item Construire deux graphes que m\^eme le GIN ne distingue pas (\'echec du 1-WL).
\end{enumerate}
\end{exercice}

\begin{exercice}[Mini-batch training]
Impl\'ementer l'entra\^inement par mini-batch avec \'echantillonnage de voisinage
pour un GCN (sans utiliser \texttt{NeighborLoader}).
\begin{enumerate}
    \item Impl\'ementer la strat\'egie d'\'echantillonnage.
    \item Comparer la vitesse d'entra\^inement avec le mode full-batch.
    \item \'Evaluer l'impact du nombre de voisins \'echantillonn\'es sur la pr\'ecision.
\end{enumerate}
\end{exercice}
